
\documentclass{article}
\usepackage{colortbl}
\usepackage{makecell}
\usepackage{multirow}
\usepackage{supertabular}

\begin{document}

\newcounter{utterance}

\centering \large Interaction Transcript for game `eqbench', experiment `eq\_bench', episode 113 with reflexion{-}llama{-}3.1{-}8b{-}instruct.
\vspace{24pt}

{ \footnotesize  \setcounter{utterance}{1}
\setlength{\tabcolsep}{0pt}
\begin{supertabular}{c@{$\;$}|p{.15\linewidth}@{}p{.15\linewidth}p{.15\linewidth}p{.15\linewidth}p{.15\linewidth}p{.15\linewidth}}
    \# & \multicolumn{2}{c}{Player} && \multicolumn{2}{c}{Game Master} \\
    \hline

    \theutterance \stepcounter{utterance}  
    & & & \multicolumn{4}{p{0.6\linewidth}}{
        \cellcolor[rgb]{0.9,0.9,0.9}{
            \makecell[{{p{\linewidth}}}]{
                \texttt{\tiny{[P1$\langle$GM]}}
                \texttt{You have attempted this same task before and did not succeed. These are the plans you wrote after those attempts.} \\
\texttt{Trial \#1: New plan: I overestimated Esther's emotional response to Ben's jabs, assuming she would feel superior due to her confident and assertive tone. However, upon re{-}reading the dialogue, I realize that Esther's responses are more defensive and dismissive, suggesting that she might be feeling a bit threatened by Ben's comments. Specifically, when Ben asks if it gets lonely being the Ice Queen, Esther responds with a somewhat harsh statement about being better off alone at the top than surrounded by fools. This suggests that she might be feeling a bit vulnerable or sensitive to Ben's words. Next time, I will consider a more nuanced emotional response, such as feeling slightly threatened or defensive, and assign a score accordingly.} \\
\\ 
\texttt{Trial \#2: New plan: I overestimated Esther's emotional response to Ben's jabs, assuming she would feel superior due to her confident and assertive tone. However, upon re{-}reading the dialogue, I realize that Esther's responses are more defensive and dismissive, suggesting that she might be feeling a bit threatened by Ben's comments. Specifically, when Ben asks if it gets lonely being the Ice Queen, Esther responds with a somewhat harsh statement about being better off alone at the top than surrounded by fools. This suggests that she might be feeling a bit vulnerable or sensitive to Ben's words. } \\
\\ 
\texttt{However, I also realize that I didn't consider the context of the dialogue and the character's personality. Esther is referred to as the "Ice Queen," which suggests that she is a confident and assertive person who is not easily intimidated. } \\
\\ 
\texttt{Considering this, I will revise my plan to take into account Esther's personality and the context of the dialogue. I will assign a higher score to "Superior" and a lower score to "Lonely" and "Scared". I will also consider the possibility that Esther might be feeling a bit annoyed or irritated by Ben's comments, which could be reflected in a higher score for "Regretful" or "Scared". However, based on the dialogue, I don't think this is the case.} \\
\\ 
\texttt{Here is my revised plan:} \\
\\ 
\texttt{Superior: 8} \\
\texttt{Lonely: 2} \\
\texttt{Regretful: 0} \\
\texttt{Scared: 0} \\
\\ 
\\ 
\\ 
\texttt{Your task is to predict the likely emotional responses of a character in this dialogue:} \\
\\ 
\texttt{Ben: Well, well, well, if it isn't the Ice Queen herself. I've heard tales of your cold{-}blooded poker face.} \\
\texttt{Esther: And you must be Benny the Bluff. I've heard you're quite the charmer. } \\
\texttt{Ben: Guilty as charged, sweetheart. But don't let that fool you. I've got a few tricks up my sleeve.} \\
\texttt{Esther: I'm sure you do. But remember, Ben, this isn't a magic show. It's a game of skill. } \\
\texttt{Ben: And a little bit of luck, wouldn't you agree?} \\
\texttt{Esther: Luck is for those who can't rely on their abilities.} \\
\texttt{Ben: Spoken like a true Ice Queen. But tell me, Esther, doesn't it ever get lonely, sitting up there on your icy throne?} \\
\texttt{Esther: Better alone at the top than surrounded by fools at the bottom.} \\
\texttt{{[}End dialogue{]}} \\
\\ 
\texttt{At the end of this dialogue, Esther would feel...} \\
\texttt{Superior} \\
\texttt{Lonely} \\
\texttt{Regretful} \\
\texttt{Scared} \\
\\ 
\texttt{Give each of these possible emotions a score from 0{-}10 for the relative intensity that they are likely to be feeling each.} \\
\\ 
\texttt{You must output in the following format, including headings (of course, you should give your own scores), with no additional commentary:} \\
\\ 
\texttt{Superior: <score>} \\
\texttt{Lonely: <score>} \\
\texttt{Regretful: <score>} \\
\texttt{Scared: <score>} \\
\\ 
\\ 
\texttt{{[}End of answer{]}} \\
\\ 
\texttt{Remember: zero is a valid score, meaning they are likely not feeling that emotion. You must score at least one emotion > 0.} \\
\\ 
\texttt{Your answer:} \\
            }
        }
    }
     \\ \\

    \theutterance \stepcounter{utterance}  
    & \multicolumn{4}{p{0.6\linewidth}}{
        \cellcolor[rgb]{0.561,0.737,0.561}{
            \makecell[{{p{\linewidth}}}]{
                \texttt{\tiny{[P1$\rangle$GM]}}
                \texttt{Superior: 9} \\
\texttt{Lonely: 0} \\
\texttt{Regretful: 0} \\
\texttt{Scared: 0} \\
            }
        }
    }
    & & \\ \\

    \theutterance \stepcounter{utterance}  
    & & & \multicolumn{2}{p{0.3\linewidth}}{
        \cellcolor[rgb]{0.9,0.9,0.9}{
            \makecell[{{p{\linewidth}}}]{
                \texttt{\tiny{[GM$|$GM]}}
                \texttt{Superior: 9} \\
\texttt{Lonely: 0} \\
\texttt{Regretful: 0} \\
\texttt{Scared: 0} \\
            }
        }
    }
    & & \\ \\

    \theutterance \stepcounter{utterance}  
    & & & \multicolumn{2}{p{0.3\linewidth}}{
        \cellcolor[rgb]{0.9,0.9,0.9}{
            \makecell[{{p{\linewidth}}}]{
                \texttt{\tiny{[GM$|$GM]}}
                \texttt{Superior: 8} \\
\texttt{Lonely: 3} \\
\texttt{Regretful: 0} \\
\texttt{Scared: 0} \\
            }
        }
    }
    & & \\ \\

    \theutterance \stepcounter{utterance}  
    & & & \multicolumn{2}{p{0.3\linewidth}}{
        \cellcolor[rgb]{0.9,0.9,0.9}{
            \makecell[{{p{\linewidth}}}]{
                \texttt{\tiny{[GM$|$GM]}}
                \texttt{game\_result = LOSE} \\
            }
        }
    }
    & & \\ \\

    \theutterance \stepcounter{utterance}  
    & & \multicolumn{4}{p{0.6\linewidth}}{
        \cellcolor[rgb]{0.95,0.95,0.95}{
            \makecell[{{p{\linewidth}}}]{
                \texttt{\tiny{[P1$\leftrightarrow$P]}}
                \texttt{Post{-}episode Reflexion (not sent to GM):} \\
\texttt{Superior: 9} \\
\texttt{Lonely: 1} \\
\texttt{Regretful: 0} \\
\texttt{Scared: 0} \\
\\ 
\texttt{New plan: I overestimated Esther's emotional response to Ben's jabs, assuming she would feel superior due to her confident and assertive tone. However, upon re{-}reading the dialogue, I realize that Esther's responses are more defensive and dismissive, suggesting that she might be feeling a bit threatened by Ben's comments. Specifically, when Ben asks if it gets lonely being the Ice Queen, Esther responds with a somewhat harsh statement about being better off alone at the top than surrounded by fools. This suggests that she might be feeling a bit vulnerable or sensitive to Ben's words.} \\
\\ 
\texttt{However, I also realize that I didn't consider the context of the dialogue and the character's personality. Esther is referred to as the "Ice Queen," which suggests that she is a confident and assertive person who is not easily intimidated.} \\
\\ 
\texttt{Considering this, I will revise my plan to take into account Esther's personality and the context of the dialogue. I will assign a higher score to "Superior" and a lower score to "Lonely" and "Scared". I will also consider the possibility that Esther might be feeling a bit annoyed or irritated by Ben's comments, which could be reflected in a higher score for "Regretful" or "Scared". However, based on the dialogue, I don't think this is the case.} \\
\\ 
\texttt{I will also pay closer attention to Esther's tone and language, as it seems to be more defensive and dismissive than I initially thought. This will help me to better understand her emotional state and assign more accurate scores.} \\
            }
        }
    }
    & \\ \\

\end{supertabular}
}

\end{document}
